\chapter{在线编程评测平台相关技术基础}

\section{自动判题机制}

在线编程评测平台最核心的任务，就是给用户提交上来的程序做自动化、标准化而且能反复进行的评价，它的基本流程一般包含接收代码、编译构建、控制运行、校验测试数据和返回结果这些步骤。跟人工批改比起来，自动评测反应更快、判定标准更统一，也更适合用在大规模的教学场景里\cite{kurnia2001online,leal2003mooshak,ala2005survey}。

自动判题在评测逻辑里的本质，就是把用户写出的程序放进受控环境运行，再对照标准输入、标准输出和提前准备好的测试用例来判定最终结果，常见结果一般有编译错误、运行错误、答案错误、超出时间限制、超出内存限制和评测通过这几种。对于做教学用途的平台来讲，自动判题不能只给出结果状态，还得尽量把运行日志、错误信息和测试记录保存下来，方便老师分析学生的提交内容，也方便平台生成更有针对性的反馈内容。

要明确一点，自动判题的评分标准必须确定、统一才行。同一份代码放在同一个运行环境，用一样的测试数据、相同的资源限制运行，必须拿到稳定不变的评分结果，这个特点就要求在线编程评测平台在做架构设计的时候，得先保证环境统一、运行隔离还有任务调度稳定可靠，不能让判题结果靠不稳定的主观分析来得到。

用在教学场景里，自动评测平台的目标也不该只停留在帮老师输出结果这一步。Ihantola 的研究提到，自动评测还能拓展到给出反馈、给学生测试能力打分还有迁移学习材料这些方向\cite{ihantola2011automated}，Tran 做过的相关项目能看出传统教学 OJ 的常见思路，就是先把题目管理、提交判题还有测试用例整理这些核心功能做完善\cite{tran2013uoj}，这些研究说明，做教学用途的评测平台，得先保证核心流程能用，之后再慢慢把反馈功能做得更强。

\begin{figure}[htbp]
  \centering
  \resizebox{0.96\textwidth}{!}{% 第2章：判题主链路流水线
\begin{tikzpicture}[
  font=\small\sffamily,
  st/.style={draw, rounded corners=2pt, minimum height=0.9cm, minimum width=2.0cm, align=center},
  arr/.style={-{Stealth}, thick}
]
  \node[st, fill=blue!8] (n1) {接收\\源代码};
  \node[st, fill=blue!8, right=0.6cm of n1] (n2) {编译\\构建};
  \node[st, fill=blue!8, right=0.6cm of n2] (n3) {受控\\运行};
  \node[st, fill=blue!8, right=0.6cm of n3] (n4) {用例\\比对};
  \node[st, fill=blue!8, right=0.6cm of n4] (n5) {结果\\落库};
  \foreach \a/\b in {n1/n2,n2/n3,n3/n4,n4/n5} \draw[arr] (\a) -- (\b);
\end{tikzpicture}
}
  \caption{在线编程评测平台判题流程示意图}
  \label{fig:ch2-judge-pipeline}
\end{figure}

\section{容器化执行与弹性支撑}

云原生不只是简单把系统放到云服务器上去运行，它是一种专门面向云环境做系统设计、开发、部署和维护的工程思路，这套思路一般会用到容器化封装、拆分服务、弹性伸缩、自动部署、资源编排和监控运行状态这些技术。对在线编程评测平台来讲，云原生技术的作用主要体现在系统架构和运行机制两个地方：一方面，这类平台需要承接大量用户同时提交代码，还要支持不同编程语言的运行环境，另一方面，判题的过程本身对资源隔离、安全管理和出问题后快速恢复，都有不低的要求，所以，在设计平台的时候加入云原生的思路，能帮助提高系统的可扩展性、稳定性还有运维的效率\cite{burns2016borg,hezhenwei2020kubernetes}。

在在线评测场景中，容器化技术首先解决的是运行环境一致性问题。学生提交的代码可能使用不同语言、不同编译器版本和不同依赖环境，如果直接在宿主机中执行，容易出现环境不一致、依赖冲突以及安全边界不清晰等问题。通过将判题任务放入容器中运行，可以预先构建统一的评测镜像，使每次代码执行都处于相对固定的环境中。同时，容器还可以配合 CPU、内存、文件系统和进程等限制手段，对用户代码的运行范围进行约束，从而降低异常程序对宿主机和其他评测任务造成影响的风险。

图~\ref{fig:ch2-container-boundary}给出了判题容器与宿主机之间的信任边界示意。

\begin{figure}[htbp]
  \centering
  \resizebox{0.92\textwidth}{!}{% 第2章：判题容器与宿主环境的隔离示意（单层 TikZ）
\begin{tikzpicture}[
  font=\small\sffamily,
  host/.style={draw, dashed, rounded corners=4pt, inner sep=12pt, minimum width=8.2cm, minimum height=3.0cm, align=center},
  ctr/.style={draw, very thick, rounded corners=3pt, fill=orange!8, inner sep=10pt, minimum width=6.2cm, align=center}
]
  \node[host] (H) {};
  \node[above=0.15cm of H.north, font=\small\bfseries] {宿主机 / 判题 Worker 节点};
  \node[ctr] at (H.center) {判题容器环境\\独立根文件系统 \quad 资源限额（CPU/内存/时间）\quad 默认禁用网络};
  \node[right=0.2cm of H.east, anchor=west, font=\scriptsize, align=left] {宿主机侧\\业务进程、数据库、\\日志与编排代理};
\end{tikzpicture}
}
  \caption{判题容器与宿主机之间的信任边界与职责划分示意}
  \label{fig:ch2-container-boundary}
\end{figure}

\section{服务解耦与异步调度}

\subsection{微服务拆分与任务队列}

微服务架构强调围绕业务能力拆分系统，让每个服务都可以独立开发、独立部署，再通过标准接口互相配合。在线编程评测平台里，用户管理、题库管理、代码提交、评测调度、结果展示和 AI 辅助反馈这些功能本身就有很清晰的职责划分，所以满足拆分服务化的基础条件\cite{balalaie2016microservices,difrancesco2019microservices}。

评测任务的调度是平台里最核心的技术难点，判题过程一般要经过编译、执行、用例校验这类比较费时间的操作，如果直接在请求线程里做同步处理，很容易引发响应堵塞，还会出现资源争抢的问题。对平台来说，采用先提交入队、再异步调度、由工作节点执行、最后回传结果这种处理模式会更合适，利用消息队列或者任务队列缓冲提交上来的任务，可以提高系统高峰时段的弹性，也能让判题资源的调度变得更清晰可控。

把微服务和任务调度结合使用后，系统可以把面向用户的请求服务、面向判题的执行服务分开处理。前者负责接收用户的提交、管理运行状态、展示最终结果，后者则负责拉取任务、运行程序、返回判题结果，这样的模式既符合在线编程评测平台的业务特点，也给多语言支持、优先级调度、节点监控这些功能打好了架构基础。

在现有的软件架构设计里，微服务拆分、限界上下文以及异步消息这些设计模式，已经经过了大量工程实践的验证，成为可以直接复用的参考模型\cite{richardson2018microservices,newman2021buildingmicroservices,fowler2014microservices,difrancesco2019microservices}。引入这些方法的意义，不是硬把整个系统拆成好多独立的代码仓库，而是让提交受理、任务排队、判题执行、结果写回还有观测审计这些不同的工作职责，在系统持续运行的过程里一直保持清晰的分界，以此降低联调的成本，缩小故障影响的范围\cite{balalaie2016microservices}。

\subsection{事件驱动与实时状态回传}

在线评测平台上，用户提交请求一般都有着比较明显的时间波动特征：平时访问量相对平稳，但到作业截止、课堂实验或者竞赛开始之后，短时间内提交量很可能会突然上涨。如果判题过程还是跟着 Web 请求链路一起同步完成，那代码编译、程序运行还有结果比对这类比较费时间的操作，就会长时间占用请求线程，很容易导致接口响应变慢、请求超时，甚至还会影响其他业务功能的正常访问，因此，更合理的处理方式是把“提交受理”和“判题执行”分开处理。系统收到用户的提交内容之后，先把提交记录保存好，给用户返回一个可以查询的提交状态，之后再交给独立的 worker 进程，异步完成编译、运行还有结果判定这些工作。这样做不光能缩短用户提交之后的等待时间，也方便把提交状态设计成排队中、编译中、运行中、已完成或者运行错误这类可以追踪的流转过程。

事件驱动与队列结合后，典型分工是：业务侧受理提交、落库与权限校验并维护对用户可见的状态；调度侧缓冲峰值、分发任务并处理重试；worker 侧完成编译运行、资源约束与结果回写。由此短请求与长任务解耦，判题节点可按负载横向扩展，调度策略迭代对前端主流程侵入较小。

异步判题通常配合状态推送：若仅靠轮询，会增加接口压力且实时性差；采用发布订阅或 WebSocket 可在状态变化时主动通知页面\cite{rfc6455websocket}。长连接需单独考虑鉴权、重连、背压与堆积等问题，否则高并发下可能成为瓶颈\cite{owasp2021top10}。

判题状态通道与编辑器语言服务均可基于 WebSocket 承载 JSON-RPC 或推送消息\cite{lsp2024spec,monacolanguageclientrepo}，但二者服务对象不同：前者面向提交结果，后者面向静态分析辅助，应在架构上分流，避免与常规业务接口混为一条链路。

\section{形成性反馈与教学支持}

传统在线编程评测平台好处是判题结果客观、流程统一标准，但只给学生返回 AC、WA 或者 CE 这类状态，没办法满足教学里给学生反馈的需求。Ihantola 研究提到，反馈的形式本身就会影响学生理解错误原因的速度\cite{ihantola2011automated}，这也就意味着，做教学用的平台，要改的地方不只是接着提高判题的速度，还要看反馈内容能不能真的帮学生接着改代码。

Ihantola 还得到另一个结论：有些自动评测的方式更适合自学的时候用，不太适合正式打分。举个例子，放在客户端执行的评测方式，虽然学习资料挪来挪去更方便，但是判出来的结果更容易被造假，所以更适合用来练习，不能用来给正式成绩打分\cite{ihantola2011automated}，这也说明，同一个平台里的自由练习、课程作业还有正式考试，不一定非得用一模一样的反馈和执行方案。

结合本文研究目标，可以把教学反馈理解为一个分层结构。底层是确定性的客观判题结果，用于保证评价一致性，中层是与编译错误、运行错误或输出差异相关的解释性信息，用于帮助学生定位问题，高层则是与代码结构、复杂度、风格和改进方向相关的形成性反馈，用于支持学习过程。本文在研究范围内不追求实现完整的视觉反馈或学生自编测试评价体系，但会将 AI 辅助模块作为中层和高层反馈的主要实现路径，以提升传统 OJ 平台在教学场景中的可用性。

图~\ref{fig:ch2-feedback-stack}给出了本文在讨论反馈设计时常用的分层示意。

\begin{figure}[htbp]
  \centering
  \resizebox{0.88\textwidth}{!}{% 第2章：反馈分层模型
\begin{tikzpicture}[
  font=\small,
  layer/.style={draw, minimum width=7.2cm, minimum height=1.0cm, align=left, inner xsep=8pt},
  arr/.style={-{Stealth}, thick}
]
  \node[layer, fill=gray!15] (L1) {\textbf{L1 客观层}：测试用例判定（AC/WA/CE/RE/TLE 等）};
  \node[layer, fill=blue!8, below=0.45cm of L1] (L2) {\textbf{L2 诊断层}：编译器/运行时错误摘要、逐用例输出片段};
  \node[layer, fill=purple!10, below=0.45cm of L2] (L3) {\textbf{L3 形成性层}：复杂度、风格、优化建议、上下文问答（概率输出）};
  \draw[arr] (L3) -- (L2);
  \draw[arr] (L2) -- (L1);
\end{tikzpicture}
}
  \caption{教学型在线评测中反馈分层示意}
  \label{fig:ch2-feedback-stack}
\end{figure}

\section{AI 辅助分析与反馈约束}

大语言模型在代码理解、错误解释、自然语言生成以及上下文推理这些方面，都表现出了不错的能力，所以它在程序设计教学里能起到一定的辅助作用。已经有研究提到，AI 模型可以结合题目描述、学生写的代码还有报错信息，生成带解释的反馈，帮学生弄明白出错的原因、找准修改的方向，还能在一定程度上补上传统在线评测系统的短板——传统系统大多只返回“通过/未通过”，或是很简短的错误提示\cite{keuning2018feedback,sarsa2022llm,jacobs2024feedback}。

但把 AI 用到在线编程评测平台的时候，不能忽略它输出内容不稳定的问题，大语言模型生成的内容，受提示语、上下文的组织方法还有模型版本的影响比较大，同一份学生提交的代码，换个条件可能就会得到完全不一样的反馈，所以，AI 更适合拿来做学习辅助工具，不适合直接参与最终给分，也没办法直接代替客观判题。而且模型生成的解释有时候说起来很顺，内容不一定完全正确；要是没有加上必要的限制，学生很可能会接纳错误的提示，最后形成不对的认知，另外教学场景还要考虑学术诚信的问题，尤其要避免模型直接给出完整的解题步骤，这样会弱化学生自己独立分析和调试代码的过程。

\section{浏览器端编辑器与语言服务器协议}

传统在线评测系统的前端大多就是题面、文本框加上提交按钮这三样，编辑器的功能一般只做到关键字着色或者简单的格式调整，对刚入门学习的人来说，这种设计最大的问题不是没有那些花里胡哨的功能，而是给出的语义反馈不够：很多错误其实在静态分析或者编译前期就能找到位置，但必须得等提交完，看完冗长的编译输出才能发现，明显拖慢了改代码的节奏。最近几年，Monaco 这类组件越来越成熟，LSP 生态也普及开来，在浏览器里用跟本地 IDE 差不多的补全、代码诊断和悬停提示变成了可以做到的事\cite{monacoeditorrepo,lsp2024spec}。

在协议设计上，LSP 把语言理解的功能从编辑器核心里拆分出来，用单独的语言服务器进程来处理代码索引、诊断还有补全计算，编辑器这边只需要负责显示内容和处理用户操作，再通过 JSON-RPC 跟服务器完成通信\cite{lsp2024spec}。对 Web 应用来说，常见的开发方式是在浏览器里调用语言客户端库，再通过 WebSocket 把 JSON-RPC 会话转接到后端或者侧车容器里的语言服务器上，比如写C/C++代码时常用的 Clangd 就是这个思路\cite{monacolanguageclientrepo}，这套方案和判题沙箱不会冲突：判题要做的是在管控好的环境里运行不可信的程序，拿到对应的测试结果，语言服务要做的是在只读或者受限的环境里分析代码，给用户提供编辑的帮助，两者在信任模型、资源分配和网络规则上分开设计就可以了。

就教学支持而言，编辑器侧语义反馈属于图~\ref{fig:ch2-feedback-stack} 所示反馈分层中的诊断层能力：它不改变客观判题结论，但能把大量低级错误前移到提交之前。语言服务器诊断与编译器错误并不完全等价：前者偏静态语义与工程化提示，后者以能否生成可执行文件为准；平台仍须保留判题链路的编译与运行反馈，而不能以 LSP 诊断替代评测结果。二者在信任模型与资源配额上相互独立；与判题链路的衔接及工程接线见第四章总体设计与第五章「Monaco 编辑器与基于 LSP 的 IDE 级语言服务」小节。

\section{本章小结}

在这一章节之中从自动判题开始一直到浏览器端 LSP 支持，包括容器执行、异步调度、实时回传、AI 反馈约束等方面，对本文实现所依赖的关键技术背景展开了梳理，还给出了与之相对应的文献、工程参照。这些内容一起成为了后续需求分析、总体设计的技术前提条件。
